%%%%%%%%%%%%%%%%%%%%%%%%
%%% 1-genesis/03.tex %%%
%%%%%%%%%%%%%%%%%%%%%%%%

\Verse{25}
{
	\hJ{וַיִּהְיוּ שְׁנֵיהֶם}
	\hJ{\cR{עֲרוּמִּ}ים}
	\hJ{הָאָדָם וְאִשְׁתּוֹ}
	\hJ{וְלֹא יִתְבֹּשָׁשׁוּ}
}
{
	\eJ{And the two of them were \cR{naked}, the human and his woman, and they were not ashamed.}
}
{

}

\Chapter{3}

\Verse{1}
{
	\hJ{וְהַנָּחָשׁ הָיָה}
	\hJ{\cR{עָרוּם}}
	\hJ{מִכֹּל חַיַּת הַשָּׂדֶה אֲשֶׁר עָשָׂה יְהֹוָה}
	\hR{אֱלֹהִים}
	\hJ{וַיֹּאמֶר אֶל הָאִשָּׁה אַף כִּי אָמַר}
	\hR{אֱלֹהִים}
	\hJ{לֹא תֹאכְלוּ מִכֹּל עֵץ הַגָּן}
}
{
	\eJ{And the snake was \cR{slier} than every animal of the field that YHWH \eR{God} had made, and he said to the woman, “Has \eR{God} indeed said you may not eat from any tree of the garden?”}
}
{
}

\Verse{2}
{
	\hJ{וַתֹּאמֶר הָאִשָּׁה אֶל הַנָּחָשׁ מִפְּרִי עֵץ הַגָּן נֹאכֵל}
}
{
	\eJ{And the woman said to the snake, “We may eat from the fruit of the trees of the garden.”}
}
{
}

\Verse{3}
{
	\hJ{וּמִפְּרִי הָעֵץ אֲשֶׁר בְּתוֹךְ הַגָּן אָמַר}
	\hR{אֱלֹהִים}
	\hJ{לֹא תֹאכְלוּ מִמֶּנּוּ וְלֹא תִגְּעוּ בּוֹ פֶּן תְּמֻתוּן}
}
{
	\eJ{But from the fruit of the tree that is within the garden \eR{God} has said, “You shall not eat from it, and you shall not touch it, or else you’ll die.”}
}
{
}

\Verse{4}
{
	\hJ{וַיֹּאמֶר הַנָּחָשׁ אֶל הָאִשָּׁה לֹא מוֹת תְּמֻתוּן}
}
{
	\eJ{And the snake said to the woman, “You won’t die!”}
}
{
}

\Verse{5}
{
	\hJ{כִּי יֹדֵעַ}
	\hR{אֱלֹהִים}
	\hJ{כִּי בְּיוֹם אֲכלְכֶם מִמֶּנּוּ וְנִפְקְחוּ עֵינֵיכֶם וִהְיִיתֶם כֵּ}
	\hR{אלֹהִים}
	\hJ{יֹדְעֵי טוֹב וָרָע}
}
{
	\eJ{Because \eR{God} knows that in the day you eat from it your eyes will be opened, and you’ll be like \eR{God}—knowing good and bad.”}
}
{
}

\Verse{6}
{
	\hJ{וַתֵּרֶא הָאִשָּׁה כִּי טוֹב הָעֵץ לְמַאֲכָל וְכִי תַאֲוָה הוּא לָעֵינַיִם וְנֶחְמָד הָעֵץ לְהַשְׂכִּיל וַתִּקַּח מִפִּרְיוֹ וַתֹּאכַל וַתִּתֵּן גַּם לְאִישָׁהּ עִמָּהּ וַיֹּאכַל}
}
{
	\eJ{And the woman saw that the tree was good for eating and that it was an attraction to the eyes, and the tree was desirable to bring about understanding, and she took some of its fruit, and she ate, and gave to her man with her as well, and he ate.}
}
{
}

\Verse{7}
{
	\hJ{וַתִּפָּקַחְנָה עֵינֵי שְׁנֵיהֶם וַיֵּדְעוּ כִּי}
	\hJ{\cR{עֵירֻמִּ}ם}
	\hJ{הֵם וַיִּתְפְּרוּ עֲלֵה תְאֵנָה וַיַּעֲשׂוּ לָהֶם חֲגֹורֹות}
}
{
	\eJ{And the eyes of the two of them were opened, and they knew that they were \cR{erom}.
	    And they picked fig leaves and made loincloths for themselves.}
}
{
    % \aA{Now they know.}

    % \aB{Know what?}

    % \aA{And they're erom, and they make clothes.}

    % \aB{Why clothes?}
    
    % \aA{Presumably to hide the erom.}
}

\Verse{8}
{
	\hJ{וַיִּשְׁמְעוּ אֶת קוֹל יְהֹוָה}
	\hR{אֱלֹהִים}
	\hJ{מִתְהַלֵּךְ בַּגָּן לְרוּחַ הַיּוֹם וַיִּתְחַבֵּא הָאָדָם וְאִשְׁתּוֹ מִפְּנֵי יְהֹוָה}
	\hR{אֱלֹהִים}
	\hJ{בְּתוֹךְ עֵץ הַגָּן}
}
{
	\eJ{And they heard the sound of YHWH \eR{God} walking in the garden in the wind of the day.
	    And the human and his woman hid from YHWH \eR{God} among the garden’s trees.}
}
{

}

\Verse{9}
{
	\hJ{וַיִּקְרָא יְהֹוָה}
	\hR{אֱלֹהִים}
	\hJ{אֶל הָאָדָם וַיֹּאמֶר לוֹ אַיֶּכָּה}
}
{
	\eJ{And YHWH \eR{God} called the human and said to him, “Where are you?”}
}
{
    % \aA{So YHWH is like ``Where are you?''}

    % \aB{Strange thing for a God to say.}

    % \aA{Who said God?}
}

\Verse{10}
{
	\hJ{וַיֹּאמֶר אֶת קֹלְךָ שָׁמַעְתִּי בַּגָּן}\\
    \hJ{וָאִ}%
	\hJ{\cR{ירָא}}
	\hJ{כִּי}
	\hJ{\cB{עֵירֹם}}
	\hJ{אָנֹכִי וָאֵחָבֵא}
}
{
	\eJ{And he said, “I heard the sound of you in the garden and was \cR{afraid}
	    because I was \cB{erom}, and I hid.”}
}
{

	% \aB{I thought they put clothes on three sentences ago.}
	
	% \aA{Who said they didn't?}
	
	% \aB{They're naked again.}

	% \aA{Who said naked?}
	
	% \aB{The man, just now, when he said he was afraid.}

	% \aA{Did he?}
	
	% \aB{You changed it.}

	% \aA{How so?}

    % \newcommand{\ghost}[1]{\makebox[\widthof{\Paleo{#1}}][l]{}}
	% \Table{|l|r|l|}{
		% \hline
		% Root & Word & Meaning \\
		% \hline
		% \Paleo{ירא} & \Paleo{\cR{\ghost{א}ירא}} & fear (v) \\
		% \hline
		% \Paleo{ירא} & \Paleo{א\cR{ירא}} & I feared \\
		% \hline
		% \Paleo{ירא} & \Paleo{י\cR{ירא}} & he will fear \\
		% \hline
		% \Paleo{ירא} & \Paleo{י\cR{ירא}ה} & fear (n) \\
		% \hline
		% \Paleo{ראה} & \Paleo{ו\cR{ירא}} & and he saw \\
		% \hline
		% \Paleo{ראה} & \Paleo{ו\cR{ירא}ו} & and they saw \\
		% \hline
	% }

}

\Verse{11}
{
	\hJ{וַיֹּאמֶר מִי הִגִּיד לְךָ כִּי עֵירֹם אָתָּה הֲמִן הָעֵץ אֲשֶׁר צִוִּיתִיךָ לְבִלְתִּי אֲכל מִמֶּנּוּ אָכָלְתָּ}
}
{
	\eJ{And He said, “Who told you that you were naked?
	    Have you eaten from the tree from which I commanded you not to eat?”}
}
{
}

\Verse{12}
{
	\hJ{וַיֹּאמֶר הָאָדָם הָאִשָּׁה אֲשֶׁר נָתַתָּה עִמָּדִי הִיא נָתְנָה לִּי מִן הָעֵץ וָאֹכֵל}
}
{
	\eJ{And the human said, “The woman, whom you placed with me, she gave me from the tree, and I ate.”}
}
{

	% \Def{Matres Lectiones}[Latin: mothers of reading]{
		% Consonant letters used to indicate vowels in Semitic writing systems.
		% In Hebrew these are primarily \heb{יהו}, and \heb{א}.
	    % These letters function both as ordinary consonants and as grammatical
	    % or vocalic markers. Their use helps disambiguate pronunciation in scripts
	    % that originally recorded mainly consonants.
	% }

}

\Verse{13}
{
	\hJ{וַיֹּאמֶר יְהֹוָה}
	\hR{אֱלֹהִים}
	\hJ{לָאִשָּׁה מַה זֹּאת עָשִׂית וַתֹּאמֶר הָאִשָּׁה הַנָּחָשׁ הִשִּׁיאַנִי וָאֹכֵל}
}
{
	\eJ{And YHWH \eR{God} said to the woman, “What is this that you’ve done?” And the woman said, “The snake tricked me, and I ate.”}
}
{
	% \Def{Matres Lectiones}[cont'd]{
		% Classical Biblical Hebrew was often written more defectively than the
		% received Masoretic text, omitting many such letters where later scribes
		% would supply them. Thus a form like \heb{פשנ} may be viewed as a defective
		% spelling of something expected to include the matres \heb{י} or \heb{ו},
		% as in \heb{פישון}. Similar phenomena occur across Hebrew, Aramaic,
		% Phoenician, Moabite, and Ugaritic orthographic traditions, especially in
		% older inscriptions.
	% }

}

\Verse{14}
{
	\hJ{וַיֹּאמֶר יְהֹוָה}
	\hR{אֱלֹהִים}
	\hJ{אֶל הַנָּחָשׁ כִּי עָשִׂיתָ זֹּאת אָרוּר אַתָּה מִכּל הַבְּהֵמָה וּמִכֹּל חַיַּת הַשָּׂדֶה עַל גְּחֹנְךָ תֵלֵךְ וְעָפָר תֹּאכַל כּל יְמֵי חַיֶּיךָ}
}
{
	\eJ{And YHWH \eR{God} said to the snake, “Because you did this, you are cursed out of every domestic animal and every animal of the field, you’ll go on your belly, and you’ll eat dust all the days of your life.”}
}
{

}

\Verse{15}
{
	\hJ{וְאֵיבָה אָשִׁית בֵּינְךָ וּבֵין הָאִשָּׁה וּבֵין זַרְעֲךָ וּבֵין זַרְעָהּ הוּא יְשׁוּפְךָ רֹאשׁ וְאַתָּה תְּשׁוּפֶנּוּ עָקֵב}
}
{
	\eJ{And I’ll put enmity between you and the woman and between your seed and her seed. He’ll strike you at the head, and you’ll strike him at the heel.”}
}
{

}

\Verse{16}
{
	\hJ{אֶל הָאִשָּׁה אָמַר הַרְבָּה אַרְבֶּה עִצְּבוֹנֵךְ וְהֵרֹנֵךְ בְּעֶצֶב תֵּלְדִי בָנִים וְאֶל אִישֵׁךְ תְּשׁוּקָתֵךְ וְהוּא יִמְשׁל בָּךְ}
}
{
	\eJ{To the woman He said, “I’ll make your suffering and your labor pain great. You’ll have children in pain. And your desire will be for your man, and he’ll dominate you.”}
}
{

}

\Verse{17}
{
	\hJ{וּלְאָדָם אָמַר כִּי שָׁמַעְתָּ לְקוֹל אִשְׁתֶּךָ וַתֹּאכַל מִן הָעֵץ אֲשֶׁר צִוִּיתִיךָ לֵאמֹר לֹא תֹאכַל מִמֶּנּוּ אֲרוּרָה הָאֲדָמָה בַּעֲבוּרֶךָ בְּעִצָּבוֹן תֹּאכְלֶנָּה כֹּל יְמֵי חַיֶּיךָ}
}
{
	\eJ{And to the human He said, “Because you listened to your woman’s voice and ate from the tree about which I commanded you saying, ‘You shall not eat from it,’ the ground is cursed on your account. You’ll eat from it with suffering all the days of your life.”}
}
{

}

\Verse{18}
{
	\hJ{וְקוֹץ וְדַרְדַּר תַּצְמִיחַ לָךְ וְאָכַלְתָּ אֶת עֵשֶׂב הַשָּׂדֶה}
}
{
	\eJ{And it will grow thorn and thistle at you, and you’ll eat the field’s vegetation.}
}
{

}

\Verse{19}
{
	\hJ{בְּזֵעַת אַפֶּיךָ תֹּאכַל לֶחֶם עַד שׁוּבְךָ אֶל הָאֲדָמָה כִּי מִמֶּנָּה לֻקָּחְתָּ כִּי עָפָר אַתָּה וְאֶל עָפָר תָּשׁוּב}
}
{
	\eJ{By the sweat of your nostrils you’ll eat bread until you go back to the ground, because you were taken from it; because you are dust and you’ll go back to dust.”}
}
{

}

\Verse{20}
{
	\hJ{וַיִּקְרָא הָאָדָם שֵׁם אִשְׁתּוֹ חַוָּה כִּי הִוא הָיְתָה אֵם כּל חָי}
}
{
	\eJ{And the human called his woman “Eve,” because she was mother of all living.}
}
{

}

\Verse{21}
{
	\hJ{וַיַּעַשׂ יְהֹוָה}
	\hR{אֱלֹהִים}
	\hJ{לְאָדָם וּלְאִשְׁתּוֹ כּתְנוֹת עוֹר וַיַּלְבִּשֵׁם}
}
{
	\eJ{And YHWH \eR{God} made skin garments for the human and his woman and dressed them.}
}
{

}

\Verse{22}
{
	\hJ{וַיֹּאמֶר יְהֹוָה}
	\hR{אֱלֹהִים}
	\hJ{הֵן הָאָדָם הָיָה כְּאַחַד מִמֶּנּוּ לָדַעַת טוֹב וָרָע וְעַתָּה פֶּן יִשְׁלַח יָדוֹ וְלָקַח גַּם מֵעֵץ הַחַיִּים וְאָכַל וָחַי לְעֹלָם}
}
{
	\eJ{And YHWH \eR{God} said, “Here, the human has become like one of us, to know good and bad. And now, in case he’ll put out his hand and take from the tree of life as well, and eat and live forever:”}
}
{

}

\Verse{23}
{
	\hJ{וַיְשַׁלְּחֵהוּ יְהֹוָה}
	\hR{אֱלֹהִים}
	\hJ{מִגַּן עֵדֶן לַעֲבֹד אֶת הָאֲדָמָה אֲשֶׁר לֻקַּח מִשָּׁם}
}
{
	\eJ{And YHWH \eR{God} put him out of the garden of Eden, to work the ground from which he was taken.}
}
{

}

\Verse{24}
{
	\hJ{וַיְגָרֶשׁ אֶת הָאָדָם וַיַּשְׁכֵּן מִקֶּדֶם לְגַן עֵדֶן אֶת הַכְּרֻבִים וְאֵת לַהַט הַחֶרֶב הַמִּתְהַפֶּכֶת לִשְׁמֹר אֶת דֶּרֶךְ עֵץ הַחַיִּים}
}
{
	\eJ{And He expelled the human, and He had the sphinxes and the flame of a revolving sword reside at the east of the garden of Eden to watch over the way to the tree of life.}
}
{

}